%% ==========================================================================
%% DRAFT -- August 11, 2026. Author byline: Daniel Kirtchakov (05oz,
%% halfounce.io; no institutional affiliation; ORCID 0009-0009-5213-4098).
%%
%% Claim stated at exactly referee-confirmed strength:
%%   * n <= 20: signed adversarial-referee verdict 2026-08-06 (429,865
%%     3-connected cubic graphs, all applicable Theorem 3.1 strong forms).
%%   * n = 22: independent referee recount COMPLETE 2026-08-11 (9/9 jobs
%%     rc=0; read 7,319,447 = A002851, conn3 5,909,292 = A204198,
%%     basefail = sfail = 0 on all 8 slices; certificate cross-check
%%     ok=5904 rejected=0 with membership + 3-connectivity).
%%   * n = 24: attack-side complete ONLY (98,101,019 graphs, zero failures,
%%     counts matching the published enumeration; the independent recount
%%     queue q_ref24.jsonl has NOT been run). Reported at that strength and
%%     no higher -- see 1.4 and 3.4. Do not inflate.
%%
%% Novelty sweep re-run on the drafting date (2026-08-11), BLOCKING, clean:
%% arXiv API (P_3-factor / 3-vertex-path abstracts; "Kelmans" in math.CO;
%% "Akiyama-Kano"), OpenAlex and Semantic Scholar on the citing literature
%% of 0910.2766, the Open Problem Garden entry fetched live, West's REGS
%% "Factors in regular graphs" page fetched live, plus general web sweeps.
%% No computational verification of the problem at any order was found.
%%
%% Framing constraints honored throughout (PROTOCOL section 11): the note
%% opens with the mathematical question, cites the computation in ONE
%% Certification paragraph (no file inventories, no chronology in the body),
%% and closes with the questions it opens.
%% ==========================================================================
\documentclass[11pt]{amsart}

\usepackage[margin=1in]{geometry}
\usepackage{amsmath,amssymb}
\usepackage{booktabs}
\usepackage{array}
\usepackage[colorlinks=true,allcolors=blue]{hyperref}

\emergencystretch=3em
\tolerance=1200
\hbadness=4000

\newcommand{\Lam}{\Lambda}
\newcommand{\vv}{v}

\theoremstyle{plain}
\newtheorem{theorem}{Theorem}[section]
\newtheorem{lemma}[theorem]{Lemma}
\newtheorem{proposition}[theorem]{Proposition}
\newtheorem{corollary}[theorem]{Corollary}
\theoremstyle{definition}
\newtheorem{definition}[theorem]{Definition}
\theoremstyle{remark}
\newtheorem{remark}[theorem]{Remark}
\newtheorem{question}[theorem]{Question}

\begin{document}

\title[Kelmans' 1984 problem verified through $22$ vertices]{Kelmans'
1984 problem on $3$-vertex path packings in cubic $3$-connected graphs:
exhaustive verification through $22$ vertices}

\author{Daniel Kirtchakov}
\address{Independent researcher (\texttt{05oz}); no institutional
affiliation. ORCID:
\href{https://orcid.org/0009-0009-5213-4098}{0009-0009-5213-4098}.}
\email{daniel@halfounce.io}
\urladdr{https://halfounce.io}

\thanks{\emph{Computation and authorship.} The searchers, the independent
verifiers, and the audits in this work were produced by \textbf{Claude}
(Anthropic), directed by the author, on a single Apple~M4 laptop. The
certificate checkers import only the Python standard library and share no
code with the searchers; the adversarial referee of \S\ref{ssec:referee}
was a separate agent instance given no access to the search pipeline's
code and wrote its own code throughout. This is a factual methods
statement, and it is part of the point of the note: the artifacts are
designed so that the provenance of the \emph{search} is irrelevant to the
validity of the \emph{result}. External tools used by the pipelines:
\texttt{geng} from \texttt{nauty}~2.9.3 \cite{nauty}, Apple clang 17,
and Python~3.}

\thanks{\emph{Prior-art record.} The primary source \cite{Kel11} was read
in full (arXiv v2; a copy is archived with the artifacts, and all
statement numbers below follow its numbering). A novelty sweep was run on
August~5 and again on August~11, 2026, the latter immediately before this
draft was written: the arXiv API (abstracts mentioning $P_3$-factors or
$3$-vertex paths; \emph{Kelmans} in math.CO; \emph{Akiyama--Kano}), the
citing literature of \cite{Kel11} through OpenAlex and Semantic Scholar,
the Open Problem Garden entry \cite{OPG} fetched live, West's REGS page
\cite{West} fetched live, and general web sweeps. The problem is recorded
as open in both problem collections, and no computational verification of
it at any order was found anywhere in the record. \emph{Draft of
August~11, 2026.}}

\subjclass[2020]{Primary 05C70; Secondary 05C40, 05C69, 68V15}
\keywords{Cubic graph, $3$-connected graph, $3$-vertex path packing,
$P_3$-factor, $\Lambda$-factor, Kelmans' problem, Akiyama--Kano
conjecture, Reed's domination conjecture, exhaustive verification,
certified computation}

\date{August 11, 2026}

\begin{abstract}
A $\Lambda$-packing of a graph is a subgraph all of whose components are
$3$-vertex paths, and a $\Lambda$-factor is a spanning one; $\lambda(G)$
denotes the maximum number of vertex-disjoint $3$-vertex paths in $G$, so
that $\lambda(G)\le\lfloor \vv(G)/3\rfloor$ always. Kelmans asked in
$1984$ whether equality holds for every cubic $3$-connected graph;
specialized to $\vv(G)\equiv 0 \bmod 6$ this is the Akiyama--Kano
conjecture that such a graph has a $P_3$-factor, and a positive answer
would give Reed's domination conjecture for cubic $3$-connected graphs.
The problem is open, and no computational verification of it at any order
appears in the record.

We supply one. For every $3$-connected cubic graph on at most $22$
vertices --- all $6{,}339{,}157$ of them --- $\lambda(G)=\lfloor
\vv(G)/3\rfloor$, and the applicable strong forms of Kelmans' equivalence
theorem hold as well: $(z2),(z3),(z7),(z8)$ at orders $6,12,18$; $(t2)$ at
$8,14,20$; $(f\mkern-2mu1),(f\mkern-2mu2)$ at $4,10,16$ and $(f\mkern-2mu1)$ at $22$. Those equivalences
are constructive, so a single failure at any of these orders would have
yielded an explicit cubic $3$-connected graph of order divisible by $6$
with no $\Lambda$-factor. None was found. The graphs were generated with
\texttt{geng}, and the generated counts agree exactly, at every order,
with the published enumerations of connected cubic graphs and of
$3$-connected cubic graphs. Two pipelines sharing no code --- different
connectivity tests, different search orders, both negative-controlled ---
agree on every count and report zero failures, and all $43{,}580$
$\Lambda$-factor certificates in the corpus were re-verified from the
\texttt{graph6} strings alone by standard-library checkers sharing no code
with the searcher that produced them ($34{,}429$ of them by two such
checkers independently). At $24$ vertices the search pipeline has
completed all $98{,}101{,}019$ $3$-connected cubic graphs with no failure
and with counts again matching the published enumeration, but the
independent recount is outstanding; that order is reported at that
strength and no higher. We close with the questions the computation
opens.
\end{abstract}

\maketitle

\section{Introduction}\label{sec:intro}

\subsection{The question}\label{ssec:question}
All graphs here are finite, undirected, and simple. Write $\vv(G)$ for
the number of vertices of $G$ and $\Lam$ for the $3$-vertex path $P_3$.
Following \cite{Kel11}, a \emph{$\Lam$-packing} of $G$ is a subgraph all
of whose components are isomorphic to $\Lam$, and a \emph{$\Lam$-factor}
is a spanning $\Lam$-packing --- what is elsewhere called a
\emph{$P_3$-factor}. Let $\lambda(G)$ be the maximum number of pairwise
disjoint $3$-vertex paths in $G$. Counting vertices gives
\begin{equation}\label{eq:trivial}
  \lambda(G)\;\le\;\left\lfloor \vv(G)/3 \right\rfloor
\end{equation}
for every graph, and the question is when \eqref{eq:trivial} is tight.

\begin{quote}
\textbf{Problem} (A. Kelmans, 1984; Problem 1.10 of \cite{Kel11}).
\emph{Is the following claim true?}
\smallskip

\noindent $(P)$ \emph{If $G$ is a $3$-connected cubic graph, then
$\lambda(G)=\lfloor \vv(G)/3\rfloor$.}
\end{quote}

\noindent A cubic graph has even order, so $(P)$ splits into three cases
by the residue of $\vv(G)$ modulo $6$: for $\vv(G)\equiv 0$ the claim is
that $G$ has a $\Lam$-factor; for $\vv(G)\equiv 4$ that some $G-x$ has
one; for $\vv(G)\equiv 2$ that some $G-\{x,y\}$ has one. The first of
these is the conjecture of Akiyama and Kano, which we take in the form
recorded by West \cite{West} --- ``when $3$ divides $n$, every
$3$-connected $3$-regular $n$-vertex graph has a $P_3$-factor'', there
attributed to the survey \cite{AK85} --- and it
is the form in which the problem is posed at the Open Problem Garden
\cite{OPG}: does every $3$-connected cubic graph on $3k$ vertices admit a
partition into $k$ paths on $3$ vertices?

Three facts fix the difficulty. First, the $\Lam$-packing problem is
$NP$-hard in general \cite{HK86} and remains $NP$-hard already for cubic
bipartite planar graphs \cite{KMZ05}, so $(P)$ is not
a statement about an easy optimization problem but about a class on which
the optimum is conjecturally forced. Second, $3$-connectivity is not a
convenience: Kelmans \cite{Kel07} constructs infinitely many $2$-connected
cubic bipartite planar graphs $G$ with $\lambda(G)<\lfloor
\vv(G)/3\rfloor$, so the hypothesis cannot be weakened by one unit of
connectivity. Third, the problem has a consequence outside path packing.
Reed \cite{Reed96} conjectured that a connected cubic graph $G$ has
domination number $\gamma(G)\le\lceil \vv(G)/3\rceil$; that conjecture is
false for connected and even for $2$-connected cubic graphs
\cite{Kel06,KS05}, but as Kelmans observes \cite[\S1]{Kel11}, if $(P)$
holds then --- by way of the strong forms of
Theorem~\ref{thm:kelmans}, not directly --- Reed's conjecture holds for
cubic $3$-connected graphs. A \emph{positive} resolution of $(P)$
therefore settles a domination question as a corollary; a counterexample
would leave Reed's conjecture on this class untouched.

The problem has been open since 1984. Partial results are known. For
every cubic graph $\lambda(G)\ge\lceil \vv(G)/4\rceil$ \cite{KM04}, and
for connected cubic graphs this has been improved to
$\lambda(G)\ge\frac{39}{152}\vv(G)$ when $\vv(G)\ge 17$ \cite{KMZ08} and
to $\lambda(G)\ge\frac{3}{11}\vv(G)$ when $\vv(G)\ge 9$ \cite{KZ08}, as
recorded in \cite[1.3, 1.6, 1.8]{Kel11}. In restricted classes the exact
statement is known: $\lambda(G)=\lfloor \vv(G)/3\rfloor$ for
$2$-connected claw-free graphs \cite{KKN01}, and claims $(z1)$--$(z5)$ of
Theorem~\ref{thm:kelmans} below hold for cubic $3$-connected claw-free
graphs \cite{Kel07b} (see also \cite{Kel11b}). But the general case is
untouched, and, as far as our sweeps can determine, no exhaustive
verification of $(P)$ at any order has ever been published.

\subsection{Result}\label{ssec:result}

\begin{theorem}\label{thm:main}
Let $G$ be a $3$-connected cubic graph with $\vv(G)\le 22$. Then
$\lambda(G)=\lfloor \vv(G)/3\rfloor$. In particular, every $3$-connected
cubic graph on $6$, $12$, or $18$ vertices has a $\Lam$-factor, i.e.\
admits a partition of its vertex set into $\vv(G)/3$ paths on $3$
vertices.
\end{theorem}

\noindent The statement ranges over all $6{,}339{,}157$ $3$-connected
cubic graphs of order at most $22$; the counts by order are in
Table~\ref{tab:counts}. Theorem~\ref{thm:main} is the base claim $(P)$.
The verification also establishes the applicable strong forms of Kelmans'
equivalence theorem, which is where its counterexample-hunting power lies
(\S\ref{ssec:strong}).

\begin{theorem}\label{thm:strong}
Let $G$ be a $3$-connected cubic graph. Then, in the ranges indicated,
\begin{itemize}
\item[$(z2)$] for $\vv(G)\in\{6,12,18\}$: $G-e$ has a $\Lam$-factor for
  every $e\in E(G)$;
\item[$(z3)$] for $\vv(G)\in\{6,12,18\}$: for every $e\in E(G)$ some
  $\Lam$-factor of $G$ contains $e$;
\item[$(z7)$] for $\vv(G)\in\{6,12,18\}$: $G-X$ has a $\Lam$-factor for
  every $X\subseteq E(G)$ with $|X|=2$;
\item[$(z8)$] for $\vv(G)\in\{6,12,18\}$: $G-L$ has a $\Lam$-factor for
  every $3$-vertex path $L$ in $G$;
\item[$(t2)$] for $\vv(G)\in\{8,14,20\}$: $G-\{x,y\}$ has a $\Lam$-factor
  for every $xy\in E(G)$;
\item[$(f\mkern-2mu1)$] for $\vv(G)\in\{4,10,16,22\}$: $G-x$ has a $\Lam$-factor
  for every $x\in V(G)$;
\item[$(f\mkern-2mu2)$] for $\vv(G)\in\{4,10,16\}$: $G-\{x,e\}$ has a $\Lam$-factor
  for every $x\in V(G)$ and every $e\in E(G)$ not incident with $x$.
\end{itemize}
\end{theorem}

\noindent Kelmans states $(f\mkern-2mu2)$ without an incidence
restriction; restricting it here to edges not incident with $x$ loses
nothing, and for the reason his own proof gives: if $e$ is incident with
$x$ then $G-\{x,e\}=G-x$, so the incident case is $(f\mkern-2mu1)$
\cite[3.21]{Kel11}.

The evidence for both theorems is a complete search: every $3$-connected
cubic graph of each order was generated, and for each the relevant
existence questions were decided exactly --- not estimated, not sampled
--- by an exhaustive procedure (\S\ref{ssec:decide}). The result was then
recounted end to end by a second pipeline written independently of the
first, and the whole was confirmed by an adversarial referee
(\S\ref{sec:verif}).

One further order was searched but is \emph{not} part of
Theorems~\ref{thm:main} and~\ref{thm:strong}, and is recorded separately
at the strength it has reached.

\begin{proposition}\label{prop:n24}
The search pipeline reports $\lambda(G)=\vv(G)/3$ --- that is, a
$\Lam$-factor --- for every one of the $98{,}101{,}019$ $3$-connected
cubic graphs on $24$ vertices, with no failure on any graph. The
enumeration underlying that sweep was audited independently (job
accounting reconstructed from the run records, generated counts matched
against the published enumerations, zero failure lines in any output),
but the graphs themselves have not been re-solved by the independent
pipeline. This order is asserted at exactly that strength.
\end{proposition}

\subsection{Why the strong forms are the counterexample hunt}
\label{ssec:strong}
Theorem~\ref{thm:strong} is not a collection of sanity checks. Its
statements are the ones Kelmans proves equivalent to $(P)$ as universal
claims over the class of cubic $3$-connected graphs.

\begin{theorem}[Kelmans {\cite[3.1]{Kel11}}]\label{thm:kelmans}
Read as universal statements over the class of cubic $3$-connected
graphs, the nineteen claims $(z1)$--$(z9)$, $(t1)$--$(t4)$,
$(f\mkern-2mu1)$--$(f\mkern-2mu6)$ --- conditioned on $\vv(G)\equiv0$,
$2$, $4 \bmod 6$ respectively, and including the seven forms listed in
Theorem~\ref{thm:strong} --- are pairwise equivalent.
\end{theorem}

\noindent Two remarks on the statement. Claim $(P)$ is not one of the
nineteen, but is equivalent to them, as Kelmans records
\cite[\S1]{Kel11}: $(z1)$ is exactly the $\vv(G)\equiv0$ case of $(P)$,
so $(P)$ implies $(z1)$, while conversely $(z1)$ yields $(t1)$ and
$(f\mkern-2mu1)$ by the theorem, and those are stronger than the
$\vv(G)\equiv 2$ and $\vv(G)\equiv 4$ cases of $(P)$. And the
quantification is over the class: the equivalences hold between claims
each universally quantified over all cubic $3$-connected graphs, not
graph by graph. A single graph failing one strong form is therefore not
by itself a graph failing another; what it yields is a construction, and
that is the next point.

The equivalences are proved by explicit graph compositions --- the
operations $A a\sigma b B$ and $Y(\cdot)$ and the graph-compositions
$B\{\cdot\}$ of \cite[\S2]{Kel11}, all of which preserve cubicity and
$3$-connectivity \cite[2.1--2.3]{Kel11} --- and Kelmans draws the
consequence himself: he gives several different proofs of $3.1$, so that
``if there is a counterexample $C$ to one of the above claims, then these
different proofs provide different constructions of counterexamples to
the other claims in 3.1'' \cite[\S1]{Kel11}. The implications therefore
run constructively in the direction that matters here: a single cubic
$3$-connected graph violating any one of $(z2),(z3),(z7),(z8),(t2),(f\mkern-2mu1)$,
or $(f\mkern-2mu2)$ can be converted, by those compositions, into an explicit cubic
$3$-connected graph of order $\equiv 0 \bmod 6$ with no $\Lam$-factor,
i.e.\ into a counterexample to $(P)$ and to the Akiyama--Kano conjecture.

This changes what the orders $\vv(G)\equiv 2,4 \bmod 6$ are worth. They
are not orders at which the factor question is vacuous; they are orders
at which a cheap local test --- delete an edge's two endpoints, or one
vertex, and ask for a factor of the remainder --- is a valid search for a
counterexample to the factor case. Of the $6{,}339{,}157$ graphs covered
by Theorem~\ref{thm:main}, only $30{,}527$ have order divisible by $3$;
the strong forms are what makes the other $99.5\%$ of the corpus
evidence rather than context.

Kelmans also records where the strong forms stop \cite[(r1)--(r8)]{Kel11},
and the boundary is close. Claim $(z7)$ is false for $|X|=3$: a $3$-edge
cut whose sides have orders not divisible by $3$ obstructs any
$\Lam$-factor of $G-X$. Claim $(z8)$ is tight for two disjoint paths ---
\cite[6.1]{Kel11} constructs cubic graphs $R_s$ with $\vv(R_s)=12s$,
cyclically $6$-connected for $s\ge2$ (the smallest member has
$c(R_1)=5$, \cite[6.1(a2)]{Kel11}), having disjoint $3$-vertex paths
$L,L'$ for which $R_s-(L\cup L')$ has no $\Lam$-factor, while $R_s-L-e$
has one for every edge $e$ and every $3$-vertex path $L$ of $R_s-e$.
Since $\vv(R_1)=12$, the exact boundary of the phenomenon lies inside the
verified range: the sweep confirms $(z8)$ for every single $3$-vertex
path at order $12$, and Kelmans' construction exhibits the two-path
failure at the same order. A verification that reported success on a strong form known to be
false would be diagnostic of a broken pipeline; none of the forms tested
here is known to be false in the tested range, and none failed.

\subsection{What is not claimed}\label{ssec:notclaimed}
We do not claim to resolve Kelmans' problem, the Akiyama--Kano
conjecture, or Reed's conjecture for cubic $3$-connected graphs. Each is
a statement about an infinite class and is untouched by any finite
computation; what a proof would have to supply is discussed in
\S\ref{sec:questions}. We do not claim a bound on the order of a putative
counterexample beyond the one Theorem~\ref{thm:main} gives, namely that
none has at most $22$ vertices. We do not claim the strong forms outside
the orders listed in Theorem~\ref{thm:strong}: in particular $(f\mkern-2mu2)$ was
not tested at order $22$, and $(z7)$ and $(z8)$ were not tested at order
$24$. We do not claim order $24$ at the strength of
Theorem~\ref{thm:main}: Proposition~\ref{prop:n24} is a report of a
completed search whose independent recount has not been performed, and it
is stated in those terms deliberately. Finally, we claim no priority over
the mathematics: the problem, the equivalence theorem, the tightness
constructions, and the link to Reed's conjecture are all Kelmans'
\cite{Kel11}, and the conjecture in its $P_3$-factor form is Akiyama and
Kano's \cite{AK85}.

\section{The verification}\label{sec:comp}

\subsection{Enumeration}\label{ssec:enum}
Each order was swept over the output of \texttt{geng} from
\texttt{nauty}~2.9.3 \cite{nauty}, invoked as
\texttt{geng -q -c -d3 -D3 $n$}, which emits one \texttt{graph6} string
per isomorphism class of connected cubic graph on $n$ vertices; the large
orders were split into residue classes by \texttt{geng}'s \texttt{res/mod}
slicing. Each pipeline then applied its own $3$-connectivity filter. No
graph was excluded on any other ground, and no symmetry reduction beyond
\texttt{geng}'s isomorph rejection was used at any stage.

The completeness of the enumeration is the load-bearing external
assumption of this note (\S\ref{sec:tcb}), and it is supported by count
agreement at every order against two independent published sources. The
number of connected cubic graphs read matches the published enumeration
of Brinkmann, Goedgebeur, and McKay \cite{BGM11} --- whose Table~1 gives
$7{,}319{,}447$ at $22$ vertices and $117{,}940{,}535$ at $24$ --- and
equivalently OEIS \texttt{A002851}. The number surviving the
$3$-connectivity filter matches the published counts of $3$-connected
cubic graphs: through $20$ vertices those of McKay and Royle
\cite{McKR86}, and through $24$ vertices OEIS \texttt{A204198}, whose
terms
\[
  0,\;1,\;2,\;4,\;14,\;57,\;341,\;2828,\;30468,\;396150,\;5909292,\;
  98101019,\;\dots
\]
(indexed by half the order, with $a(11)$ and $a(12)$ credited to Ed Wynn,
2023) reproduce our filtered counts exactly at every order. The two
checks are independent of each other in the useful direction: the first
tests the generator, the second tests the filter.

\begin{table}[t]
\centering
\begin{tabular}{r r r l l}
\toprule
$n$ & connected cubic & $3$-connected & strong forms verified & recount \\
\midrule
$4$  & $1$ & $1$ & $(f\mkern-2mu1),(f\mkern-2mu2)$ & yes \\
$6$  & $2$ & $2$ & $(z2),(z3),(z7),(z8)$ & yes \\
$8$  & $5$ & $4$ & $(t2)$ & yes \\
$10$ & $19$ & $14$ & $(f\mkern-2mu1),(f\mkern-2mu2)$ & yes \\
$12$ & $85$ & $57$ & $(z2),(z3),(z7),(z8)$ & yes \\
$14$ & $509$ & $341$ & $(t2)$ & yes \\
$16$ & $4{,}060$ & $2{,}828$ & $(f\mkern-2mu1),(f\mkern-2mu2)$ & yes \\
$18$ & $41{,}301$ & $30{,}468$ & $(z2),(z3),(z7),(z8)$ & yes \\
$20$ & $510{,}489$ & $396{,}150$ & $(t2)$ & yes \\
$22$ & $7{,}319{,}447$ & $5{,}909{,}292$ & $(f\mkern-2mu1)$ & yes \\
\midrule
$\le 22$ & $7{,}875{,}918$ & $6{,}339{,}157$ & & \\
\midrule
$24$ & $117{,}940{,}535$ & $98{,}101{,}019$ & --- (base claim only) &
  \emph{pending} \\
\bottomrule
\end{tabular}
\medskip
\caption{Coverage by order. The ``connected cubic'' column is the
generator's output and matches \cite{BGM11} and OEIS \texttt{A002851};
the ``$3$-connected'' column is what survives each pipeline's own filter
and matches OEIS \texttt{A204198} at every order and \cite{McKR86} at
orders $10$ through $20$. Every graph counted
in the third column was decided for the base claim $(P)$ and for the
strong forms listed. The last row is Proposition~\ref{prop:n24}: complete
on the search side, not yet recounted independently.}
\label{tab:counts}
\end{table}

\subsection{Deciding each graph}\label{ssec:decide}
For a cubic graph on $n\le 31$ vertices the whole question fits in
machine words, and both pipelines decide it exactly rather than
approximately. The base claim asks for a set $S$ of $n \bmod 3$ vertices
such that $G-S$ has a $\Lam$-factor; a $\Lam$-factor is found, or proved
absent, by depth-first search over vertex triples $(a,b,c)$ with
$ab,bc\in E(G)$, extending a partial packing until the covered set is
everything, with a hash cache of covered-set masks already shown to
admit no completion. The strong forms are decided by running the same
factor decision on each of the linearly or quadratically many derived
graphs: $G-e$ and $G$-with-$e$-forced for $(z2),(z3)$; $G-\{e,f\}$ over
all $\tbinom{|E|}{2}$ pairs for $(z7)$; $G-V(L)$ over all $3$-vertex paths
for $(z8)$; $G-\{x,y\}$ over all edges for $(t2)$; $G-x$ over all
vertices for $(f\mkern-2mu1)$; $G-\{x,e\}$ for $(f\mkern-2mu2)$. Because the procedure is
exhaustive, a negative answer is a proof of non-existence, which is why a
single failure anywhere would have been the object of interest rather
than a nuisance.

The second pipeline decides the same questions by deliberately different
means: connectivity by union--find rather than bitmask breadth-first
search; $3$-vertex-path triples precomputed into per-vertex tables rather
than generated on the fly; depth-first branching on the highest uncovered
vertex rather than the lowest; the base claim by explicit iteration over
avoided sets rather than by a skip counter; and a four-way probed failure
cache rather than a single-slot one. It shares no code with the first.
Agreement between the two is therefore agreement between two search
orders and two data structures, not between two runs of one program.

\subsection{Certificates}\label{ssec:certs}
A completed packing is a short, checkable object, and the sweeps emit it:
one line per certified graph carrying the \texttt{graph6} string, the
list of vertex triples forming the paths, and the avoided vertices, for
example
\begin{quote}\small\ttfamily\raggedright
CERT U???????C?W?[?Y?C`Cc?Aa?X??BG?I\_?Ao?@K?? | 10-0-11 12-1-13 2-14-5
15-3-17 4-16-6 18-7-19 20-8-21 | 9
\end{quote}
at order $22$, where the seven triples and the single avoided vertex $9$
partition the $22$ vertices. Certificates were stored for every
$3$-connected cubic graph through order $18$ and for a fixed sample at
higher orders ($43{,}580$ lines through order $22$; $9{,}776$ more at
order $24$). A checker re-derives everything from the \texttt{graph6}
string alone: it decodes the string with its own decoder, verifies that
the graph is simple and cubic, optionally verifies membership of the
string in a freshly generated \texttt{geng} stream for that order and
re-proves $3$-connectivity by exhaustive deletion of vertex pairs, checks
that each listed triple is a path of the graph, and checks that the
triples and the avoided vertices partition the vertex set with
$|{\rm avoided}| = n \bmod 3$. Certificates are not the primary evidence
--- full coverage rests on the searches themselves --- but they make
spot-checking cheap for a reader, and they are regenerable: the searchers
are deterministic, so any certificate reproduces from the recorded
command.

\section{Independent recount, controls, and referee}\label{sec:verif}

\subsection{The recount}\label{ssec:recount}
The second pipeline re-ran the sweep from the generator forward. Through
order $16$ it re-solved the full stream with all strong forms enabled;
at order $18$ with $(z2),(z3),(z7),(z8)$; at order $20$ with $(t2)$; at
order $22$ with $(f\mkern-2mu1)$, in eight generator slices. At every order the
count of graphs read and the count surviving the $3$-connectivity filter
agree with the first pipeline and with the published enumerations, and
the number of base-claim failures and strong-form failures is zero on
every slice. At order $22$ the eight slices read
$518{,}580 + 670{,}736 + 982{,}556 + 1{,}068{,}280 + 1{,}569{,}040 +
879{,}959 + 978{,}267 + 652{,}029 = 7{,}319{,}447$ graphs and kept
$5{,}909{,}292$, both exactly the published values; the order-$22$
recount and its certificate cross-check together consumed under two hours
of single-core time on the laptop.

\subsection{Certificate cross-checks}\label{ssec:certcheck}
The certificate corpus was re-verified across the pipeline boundary. All
$30{,}468$ certificates at order $18$, all $3{,}961$ at order $20$, and
all $5{,}904$ at order $22$ --- $40{,}333$ in total --- were accepted by
the second pipeline's own checker, with zero rejections; at orders $20$
and $22$ the check included re-proving $3$-connectivity from the
\texttt{graph6} string, and at all three orders it included membership in
a freshly regenerated generator stream ($41{,}301$, $510{,}489$, and
$7{,}319{,}447$ lines respectively, each line count re-derived at check
time). The $3{,}247$ certificates at orders $\le 16$ were accepted by the
first pipeline's independent standard-library checker, and those graphs
were in addition re-solved from scratch by the second pipeline.

\subsection{Negative controls}\label{ssec:negctl}
A sweep that only ever reports success proves nothing about its own
failure path, so both tools were forced to fail. With the
$3$-connectivity filter disabled, the base sweep over all \emph{connected}
cubic graphs of orders $10$ through $16$ reports exactly one failure: the
$16$-vertex graph with \texttt{graph6} string
\texttt{O???E?oBEAWOKGK\_@o?W\_}, which has a cut vertex. The same graph
had been identified independently by the first pipeline's own control,
and a third, unrelated brute-force implementation confirms $\lambda=4<5$
for it; the two pipelines agree on the unique sub-$3$-connected failure in
that range. With the filter disabled the strong-form paths fire as well:
no failure at order $8$, where none exists; $54$ at order $10$ ($2$ of
type $(f\mkern-2mu1)$, $52$ of $(f\mkern-2mu2)$); $253$ at order $12$,
spread over all four $(z\cdot)$ paths ($4$ of $(z2)$, $8$ of $(z3)$, $106$
of $(z7)$, $135$ of $(z8)$); $145$ of $(t2)$ type at order $14$; and
$15{,}691$ at order $16$ ($317$ of $(f\mkern-2mu1)$, $15{,}374$ of
$(f\mkern-2mu2)$). Each checker was likewise run on deliberately doctored
certificates, one corruption per line: a non-path triple, a cover gap, an
overlapping triple, a valid partition of sub-maximum shape, a
\texttt{graph6} string with a trailing character, a \texttt{graph6}
string with a flipped character, a valid certificate for a non-canonical
relabelling of the graph (a string the generator never emits), and a
genuine $\Lam$-factor certificate for a connected cubic graph that is not
$3$-connected. Every one is rejected, each by the gate it targets and
with a nonzero exit code, while the intact certificate in the same file
is accepted; and two controls on the controls confirm the diagnosis, since
disabling the membership test admits the relabelling and nothing else,
and disabling the $3$-connectivity test admits the sub-$3$-connected
certificate and nothing else. All controls in this subsection were re-run
on the drafting date against the current binaries and checkers.

\subsection{The adversarial referee}\label{ssec:referee}
The result rests on an adversarial referee pass by a separate agent
instance under an explicit independence rule: every load-bearing check
used code written in that session, sharing nothing with the search
pipeline beyond the \texttt{graph6} format specification. The referee
wrote the second pipeline, ran the recount and the cross-checks and the
controls described above, recomputed every count sum mechanically from
the raw per-slice summary lines rather than accepting reported totals,
and fetched the published enumeration values itself. On this protocol the
range $\vv(G)\le 20$ was confirmed on August~6, 2026, and order $22$ was
recounted to completion on August~11, 2026; Theorems~\ref{thm:main}
and~\ref{thm:strong} are stated at exactly the strength those passes
support. Order $24$ was not recounted. What the referee did establish
there is the completeness of the search side --- that every job finished
successfully, that the generated and filtered counts equal the published
enumeration values, and that no output file anywhere contains a failure
line --- which is why Proposition~\ref{prop:n24} is phrased as a report
of a completed search rather than as a verification.

\section{Certification}\label{sec:certification}

The permanent record of the computation is the part a reader can check
without us: the two certificate checkers in source form, the recorded
command line for every sweep and every generator slice at every order,
the per-slice summary outputs, the full certificate corpus, the
negative-control inputs and outputs, and the referee's verdict records.
It ships with this note and will be deposited on publication in the
program's public certificate repository \emph{Certify}
(\url{https://github.com/05oz/certify}; concept DOI
\href{https://doi.org/10.5281/zenodo.21799111}{10.5281/zenodo.21799111}).
The two searchers are specified in \S\ref{ssec:decide} but are not part
of that deposit, and we say so rather than imply otherwise: what is
deposited is what re-checks their output without them, which is every
positive answer they gave. Nothing large needs to be archived either --- the
graph streams are outputs of \texttt{geng} and regenerate from the
recorded commands. To replay: checking any certificate, or the whole
corpus, requires only CPython; regenerating any graph stream requires
\texttt{nauty}~2.9.3. The full sweep through order $22$ is about an hour
and a half of single-core time per pipeline; order $24$ was about $29$
hours of completing-run time on the search side, and about $48$ hours of
machine time once the slices that exceeded their first time limit and
were rerun are counted.

\section{The trusted base}\label{sec:tcb}

What a skeptic must believe, separated from what they can replay.

\emph{Machine-checked, replayable:} the decision, for each of the
$6{,}339{,}157$ graphs, of the base claim and of the strong forms listed
in Theorem~\ref{thm:strong}, by each of two pipelines sharing no code;
the count of graphs read and of graphs kept at every order, by both
pipelines and against two published enumerations; all $43{,}580$
certificates through order $22$, re-verified from the \texttt{graph6}
strings alone, $34{,}429$ of them by two unrelated checkers; and the
negative controls, which both pipelines' failure paths pass.

\emph{Trusted:}
\begin{enumerate}
\item \emph{Completeness of the enumeration.} That
\texttt{geng -c -d3 -D3 $n$} emits at least one representative of every
isomorphism class of connected cubic graph on $n$ vertices is assumed,
not proved here. It is the only assumption a counterexample could hide
behind, and the mitigation is count agreement at every order with the
independent published enumerations of \cite{BGM11} and \cite{McKR86} ---
values obtained by different generators (\texttt{minibaum},
\texttt{snarkhunter}, and McKay and Royle's constructions) --- and, for
the filtered counts, with OEIS \texttt{A204198}. A generator that missed
a class would have to miss it while still producing the published totals.
\item \emph{The $3$-connectivity filters.} Both pipelines test
$3$-connectivity directly, by verifying that $G$, every $G-u$, and every
$G-\{u,v\}$ is connected, rather than by any equivalence between vertex
and edge connectivity for cubic graphs. The two implementations differ
(breadth-first search on bitmasks; union--find) and agree on every count;
and the standard-library checker's own $3$-connectivity routine, driven
over the raw generator stream by a separate driver, reproduced the
filtered counts at orders $8$ through $16$, closing the question in both
directions at those orders.
\item \emph{The correctness of the exhaustive search.} A ``no'' from
either searcher is a claim of non-existence and is not certifiable by a
short witness. Nothing in Theorems~\ref{thm:main} and~\ref{thm:strong}
depends on such a ``no'' being right --- every reported outcome is a
``yes'', witnessed by a packing --- but the \emph{absence of failures}
does depend on the searchers not silently skipping graphs, which is why
the counts, and not only the verdicts, are cross-checked at every order.
\item \emph{The equivalences of Theorem~\ref{thm:kelmans}}, on which the
counterexample-hunting reading of the strong forms rests. These are
Kelmans' theorems, proved on paper in \cite{Kel11}; they are not
machine-checked here.
\item \emph{CPython, the C compiler, the operating system, and the
hardware.}
\end{enumerate}

\noindent Neither searcher is in the trusted base in the usual sense:
every positive answer carries a packing that a reader can re-check
without trusting the program that found it.

\section{Questions}\label{sec:questions}

\subsection{What a structural proof would need}\label{ssec:qproof}
The computation supplies no tightness anywhere. Every one of the
$6{,}339{,}157$ graphs meets the bound \eqref{eq:trivial}, and meets it
after the deletions demanded by the strong forms as well; nothing in the
range distinguishes hard instances from easy ones. That is evidence for
$(P)$ and simultaneously evidence that a proof will not come from
extremal considerations at small orders.

Theorem~\ref{thm:kelmans} says a proof needs only one of the equivalent
claims, and the local ones look most tractable: $(t2)$ asserts that
$G-\{x,y\}$ has a $\Lam$-factor for \emph{every} edge $xy$, and $(f\mkern-2mu1)$
that $G-x$ has one for \emph{every} vertex --- universally quantified
statements about a single deletion, the natural shape for a discharging
or a minimal-counterexample argument. The known obstructions say where
such an argument must be careful: by \cite[(r1)]{Kel11} a $3$-edge cut
with both sides of order not divisible by $3$ kills $\Lam$-factors of
$G-X$, and by \cite[(r4)]{Kel11} claim $(t2)$ is false for non-adjacent
$x,y$. Adjacency and the divisibility of the two sides of a cut are
exactly the hypotheses that cannot be dropped.

\begin{question}\label{q:cyclic}
Does $(P)$ reduce to cyclically $4$- or $5$-connected cubic graphs? That
is, do the compositions of \cite[\S2]{Kel11} let a minimum
counterexample be assumed to have large cyclic connectivity, and if so,
does the $\Lam$-factor problem become tractable on that class --- where,
by \cite[6.1]{Kel11}, cyclically $6$-connected examples with delicate
behaviour already exist, from $24$ vertices upward?
\end{question}

\begin{question}\label{q:clawfree}
Kelmans proves $(z1)$--$(z5)$ for cubic $3$-connected \emph{claw-free}
graphs \cite{Kel07b,Kel11b}. The class is narrow: in a cubic claw-free graph the
three neighbours of any vertex cannot be pairwise non-adjacent, so every
vertex lies in a triangle. What is the weakest local-structure hypothesis
--- a bound on the number of independent claws, say --- under which the
claw-free argument still runs?
\end{question}

\subsection{The strong forms' status}\label{ssec:qstrong}
Theorem~\ref{thm:strong} is complete for the forms it lists and silent
elsewhere. Three gaps are immediate. Claim $(f\mkern-2mu2)$ was tested only through
order $16$; at order $22$ its cost is the number of vertex--edge pairs
times a factor decision, an order of magnitude above $(f\mkern-2mu1)$, and it was
not run. Claims $(z7)$ and $(z8)$ have never been tested at order $24$,
which is the first order divisible by $6$ beyond the verified range and
therefore the first place where a $(z\cdot)$ failure would be a direct
counterexample rather than a converted one. And the remaining forms of
Theorem~\ref{thm:kelmans} --- $(z4),(z5),(z6),(z9),(t1),(t3),(t4)$,
$(f\mkern-2mu3)$--$(f\mkern-2mu6)$ --- have not been tested at any order. They are equivalent
to $(P)$, so testing them adds no logical strength, but they are
different searches with different failure modes, and one of them may be
cheap enough to push several orders past $22$ on its own. Identifying
which is a concrete question.

\begin{question}\label{q:cheapest}
Among the claims of Theorem~\ref{thm:kelmans}, which admits the cheapest
exhaustive test per graph at order $n$? A form testable in time
comparable to a single factor decision would move the frontier further
than any improvement to the base search.
\end{question}

\subsection{Where the frontier sits}\label{ssec:qfrontier}
The binding constraint is enumeration, not search. Order $24$ took about
$29$ hours of completing-run single-core time on the search side for
$117{,}940{,}535$ generated graphs, a throughput near a thousand graphs
per second; the outstanding work at that order is the independent
recount, which is a comparable cost and would raise
Proposition~\ref{prop:n24} to the strength of Theorem~\ref{thm:main}.
Order $26$ requires generating $2{,}094{,}480{,}864$ connected cubic
graphs and deciding $1{,}782{,}392{,}646$ of them --- roughly three weeks
of single-core time per pipeline at the order-$24$ throughput, and more
than that in truth, since per-graph cost grows with order; it is
embarrassingly parallel across generator slices and so is a matter of
core-days, not of algorithms. Order $26$ is $\equiv 2 \bmod 6$, so it
tests $(t2)$; order $28$ ($\equiv 4$) tests $(f\mkern-2mu1)$ and has
$35{,}085{,}504{,}243$ graphs to decide.

The real wall is at order $30$, the next order divisible by $6$ after
$24$. OEIS \texttt{A204198} currently stops at $28$ vertices, so beyond
that order there is no published count of $3$-connected cubic graphs to
check a generator against --- and count agreement with an independent
enumeration is the mitigation that makes the trusted base of
\S\ref{sec:tcb} tolerable. Past $28$ vertices a sweep would be asserting
its own completeness.

\begin{question}\label{q:restricted}
Can the frontier be advanced much further on a restricted class? The
obstructions of \cite[(r1)--(r4)]{Kel11} are cuts and short cycles, which
suggests looking where neither is available; but the suggestion is not
safe, since \cite[(r5), (r6)]{Kel11} exhibit obstructions to the
$(f\mkern-2mu\cdot)$ forms in cubic $3$-connected graphs with no
$3$-cycles and no $4$-cycles. Generators for sparse families ---
\texttt{genreg} \cite{genreg},
\texttt{snarkhunter} \cite{BGM11} --- reach far higher orders than the
full cubic enumeration does, because the families are far sparser.
A sweep of, say, all cyclically $5$-connected cubic graphs of order
$\equiv 0 \bmod 6$ up to $34$ or $36$ vertices looks feasible today, and
would test the conjecture on the class where the compositions of
\cite[\S2]{Kel11} have the least room to work.
\end{question}

\section*{Acknowledgments}

The problem is Alexander Kelmans'; the equivalence theorem, the tightness
constructions, and the connection to Reed's domination conjecture that
give this verification its meaning are all his, and this note is a
computation performed inside his framework. The $P_3$-factor form of the
conjecture is due to Jun-ichi Akiyama and Mikio Kano, and its continued
visibility owes much to Douglas West's problem pages and to the Open
Problem Garden. The enumeration rests on Brendan McKay's \texttt{nauty},
and the count cross-checks on the published cubic-graph enumerations of
Brinkmann, Goedgebeur, and McKay and of McKay and Royle, and on the OEIS.
The computation and drafting were AI-assisted as stated in the first
footnote; the adversarial referee protocol of \S\ref{ssec:referee}, its
rule that a claim be stated at exactly the strength its recount supports,
and its rule that a computation-only note close with the questions it
opens shaped this note's final form.

\begin{thebibliography}{99}
\footnotesize

\bibitem[AK85]{AK85}
J.~Akiyama and M.~Kano,
\emph{Factors and factorizations of graphs --- a survey},
J.\ Graph Theory \textbf{9} (1985), no.~1, 1--42. Updated and augmented
in: \emph{Factors and Factorizations of Graphs}, Lecture Notes in Math.\
\textbf{2031}, Springer, 2011.

\bibitem[BGM11]{BGM11}
G.~Brinkmann, J.~Goedgebeur and B.~D. McKay,
\emph{Generation of cubic graphs},
Discrete Math.\ Theor.\ Comput.\ Sci.\ \textbf{13} (2011), no.~2,
69--80. Table~1 supplies the connected cubic counts used here; the
generator is \texttt{snarkhunter}.

\bibitem[HK86]{HK86}
P.~Hell and D.~G. Kirkpatrick,
\emph{Packings by complete bipartite graphs},
SIAM J.\ Algebraic Discrete Methods \textbf{7} (1986), no.~2, 199--209.

\bibitem[KKN01]{KKN01}
A.~Kaneko, A.~Kelmans and T.~Nishimura,
\emph{On packing $3$-vertex paths in a graph},
J.\ Graph Theory \textbf{36} (2001), 175--197.

\bibitem[Kel06]{Kel06}
A.~Kelmans,
\emph{Counterexamples to the cubic graph domination conjecture},
arXiv:math/0607512 (2006).

\bibitem[Kel07]{Kel07}
A.~Kelmans,
\emph{Packing $3$-vertex paths in $2$-connected graphs},
arXiv:0712.4151 (2007); also RUTCOR Research Report RRR 21--2005,
Rutgers University (2005).

\bibitem[Kel07b]{Kel07b}
A.~Kelmans,
\emph{Packing $3$-vertex paths in claw-free graphs},
arXiv:0711.3871 (2007).

\bibitem[Kel11]{Kel11}
A.~Kelmans,
\emph{Packing $3$-vertex paths in cubic $3$-connected graphs},
arXiv:0910.2766v2 (July~25, 2011); the text records acceptance by
Discrete Math.\ on August~14, 2009, but no journal version has appeared,
and the arXiv text is the version of record used here. Problem~1.10 is
the 1984 problem; Theorem~3.1 is the equivalence theorem; \S6 is the
$R_s$ construction.

\bibitem[Kel11b]{Kel11b}
A.~Kelmans,
\emph{Packing $3$-vertex paths in claw-free graphs and related topics},
Discrete Appl.\ Math.\ \textbf{159} (2011), no.~2--3, 112--127.

\bibitem[KM04]{KM04}
A.~Kelmans and D.~Mubayi,
\emph{How many disjoint $2$-edge paths must a cubic graph have?},
J.\ Graph Theory \textbf{45} (2004), 57--79.

\bibitem[KMZ05]{KMZ05}
A.~Kosowski, M.~Ma\l afiejski and P.~\.Zyli\'nski,
\emph{Parallel processing subsystems with redundancy in a distributed
environment},
in: Parallel Processing and Applied Mathematics (PPAM 2005), Lecture
Notes in Comput.\ Sci.\ \textbf{3911}, Springer, 2006, 1002--1009; cited
in \cite{Kel11} for the $NP$-hardness of $3$-vertex path packing on cubic
bipartite planar graphs.

\bibitem[KMZ08]{KMZ08}
A.~Kosowski, M.~Ma\l afiejski and P.~\.Zyli\'nski,
\emph{Tighter bounds on the size of a maximum $P_3$-matching in a cubic
graph},
Graphs Combin.\ \textbf{24} (2008), no.~5, 461--468.

\bibitem[KS05]{KS05}
A.~V. Kostochka and B.~Y. Stodolsky,
\emph{On domination in connected cubic graphs},
Discrete Math.\ \textbf{304} (2005), no.~1--3, 45--50.

\bibitem[KZ08]{KZ08}
A.~Kosowski and P.~\.Zyli\'nski,
\emph{Packing three-vertex paths in $2$-connected cubic graphs},
Ars Combin.\ \textbf{89} (2008). (The bound quoted here is the one
recorded as \cite[1.8]{Kel11}.)

\bibitem[McKR86]{McKR86}
B.~D. McKay and G.~F. Royle,
\emph{Constructing the cubic graphs on up to $20$ vertices},
Ars Combin.\ \textbf{21A} (1986), 129--140. Author's copy:
\url{https://users.cecs.anu.edu.au/~bdm/papers/Gobstoppers.pdf}; the
connectivity tables there give the $3$-connected counts used above.

\bibitem[McKP14]{nauty}
B.~D. McKay and A.~Piperno,
\emph{Practical graph isomorphism, II},
J.\ Symbolic Comput.\ \textbf{60} (2014), 94--112. Version used:
\texttt{nauty}~2.9.3 (\texttt{geng}).

\bibitem[Mer99]{genreg}
M.~Meringer,
\emph{Fast generation of regular graphs and construction of cages},
J.\ Graph Theory \textbf{30} (1999), 137--146. (\texttt{genreg}.)

\bibitem[OPG]{OPG}
\emph{Partition of a cubic $3$-connected graph into paths of length $2$},
Open Problem Garden, posed by A.~Kelmans, posted March~4, 2013,
\url{http://www.openproblemgarden.org/op/partition_of_a_cubic_3_connected_graphs_into_paths_of_length_2},
accessed 2026-08-11.

\bibitem[Reed96]{Reed96}
B.~Reed,
\emph{Paths, stars and the number three},
Combin.\ Probab.\ Comput.\ \textbf{5} (1996), 277--295.

\bibitem[West]{West}
D.~B. West,
\emph{Factors in regular graphs}, REGS problem page,
\url{https://dwest.web.illinois.edu/regs/facreg.html}, accessed
2026-08-11.

\end{thebibliography}

\end{document}
